
\section{每个光滑三次曲面（$\bar{k}$ 上）都有 27 条线}

我们现已准备好证明\mbox{定理~27.0.1}。在到达\mbox{练习~27.2.K}之前，为免分心，我们假设 $\operatorname{char}\bar{k}=0$。不过，即便 $\operatorname{char}\bar{k}\neq 3$，下面的论证也无须改动即可通过。唯一需要的一处检验——也正是对特征加以限制的原因——是\mbox{练习~27.1.C} 在把 $\mathbb{C}$ 换成任何一个这样的 $\bar{k}$ 之后依然成立。

\textbf{27.2.A.（练习）}\quad（两小问共用提示：回想\mbox{练习~12.2.N} 的解答。）
（a）定义\trIncidenceCorrespondence{} $Z\subset\mathbb{P}^{19}\times\mathbb{G}(1,3)$，它对应于「$\mathbb{P}^{3}$ 中含于某张三次曲面 $X$ 的一条直线 $\ell$」这份数据。（与\mbox{练习~27.1.D} 一样，这也是练习的一部分。）令 $\pi$ 为\trProjection{} $Z\to\mathbb{P}^{19}$。
\[
\begin{tikzcd}[row sep=2.2em, column sep=3.6em]
& Z \ar[dl, "\pi"] \ar[dr, "\mathbb{P}^{15}\text{-丛}"] & \\
\mathbb{P}^{19} && \mathbb{G}(1,3)
\end{tikzcd}
\]
（b）通过把 $Z$ 描述为\trGrassmannian{} $\mathbb{G}(1,3)$ 上的一个 $\mathbb{P}^{15}$-丛，证明 $Z$ 是维数为 $19$ 的不可约光滑\trVariety。

\textbf{27.2.B.（练习）}\quad 利用「存在只有有限多条直线的三次曲面」（\mbox{练习~27.1.C}）这一事实，以及态射纤维维数的行为（\mbox{练习~12.4.A} 或\mbox{定理~12.4.1}），证明下述诸条。
（a）每张三次曲面都含有直线，也就是说 $\pi$ 是满射。（提示：证明 $\pi$ 是\trProjective{}的。）
（b）「大多数三次曲面只有有限多条直线」：存在一个\trDense{}\trOpenSubset{} $U\subset\mathbb{P}^{19}$，使得由 $U$ 的\trClosedPoint{}所参数化的那些三次曲面都有条数为正的有限多条直线。（提示：纤维维数的\trUpperSemicontinuous{}性，\mbox{定理~12.4.3}。）

下面这个事实是\mbox{定理~27.0.1} 证明中的关键结果，也是 27 条线的主要奇迹之一：它保证在一张光滑曲面上，这些直线能保持两两相异。它的非正式陈述是：两条线不可能凑到一起而不损伤曲面。这其实是一条\trDeformationTheory{}的结果：我们正在显式地证明，光滑三次曲面中的一条直线没有任何一阶形变。

\textbf{27.2.1. \trhTheorem。}\ ——\ \emph{若 $\ell$ 是一张正则三次曲面 $X$ 中的一条直线，则 $\{\ell\subset X\}$ 是 $\pi\colon Z\to\mathbb{P}^{19}$ 的纤维中的一个\trReduced{}孤立点。}

在证明\mbox{定理~27.2.1} 之前，我们先用它来证明\mbox{定理~27.0.1}。

\textbf{27.2.2.}\ \emph{\mbox{定理~27.0.1} 的证明（假定\mbox{定理~27.2.1}）。}\quad 此时 $\pi$ 是一个\trProjective{}\trMorphism，而由\mbox{定理~27.2.1}，在 $\mathbb{P}^{19}\setminus\Delta$ 之上，$\pi$ 的相对维数为 $0$，从而具有有限的\trFiber。于是，由\mbox{定理~18.1.6}（射影＋有限纤维＝有限），$\pi$ 在 $\mathbb{P}^{19}\setminus\Delta$ 上是有限的。

此外，$Z$ 与 $\mathbb{P}^{19}$ 都（在 $\mathbb{Z}$ 上）光滑且维数为 $19$，而 $\pi\colon Z\to\mathbb{P}^{19}$ 在 $\mathbb{P}^{19}\setminus\Delta$ 上方的\trFiber{}是（几何）既约的点，这仍由\mbox{定理~27.2.1} 得出。再由\mbox{练习~24.8.G}，$\pi$ 在 $\mathbb{P}^{19}\setminus\Delta$ 上是\trEtale{}的（并且，特别地，是\trFlat{}的）。

于是，在 $\mathbb{P}^{19}\setminus\Delta$ 之上，$\pi$ 是一个有限\trFlat{}\trMorphism，从而 $\pi$ 的\trFiber（再次强调，在远离 $\Delta$ 处）总有相同个数的点——「按正确方式计数」（\mbox{注记~24.3.8}）。但由\mbox{定理~27.2.1}，在 $\mathbb{P}^{19}\setminus\Delta$ 的每个\trClosedPoint{}上方，纤维的每个点都以重数 $1$ 计入（这里用到 $k=\bar{k}$）。最后，由\mbox{练习~27.1.C}，\trFermat{}三次曲面给出了一张恰有 27 条直线的正则三次曲面的例子，所以（由于 $\mathbb{P}^{19}\setminus\Delta$ 是\trConnected{}的）我们便证毕了。\hfill$\square$

我们实际上已经证明了：在远离 $\Delta$ 处，$Z\to\mathbb{P}^{19}$ 是一个次数为 $27$ 的有限平展态射。

\textbf{27.2.3.}\ \hard\ \emph{\mbox{定理~27.2.1} 的证明。}\quad 选取射影坐标，使得直线 $\ell$ 在一个特异仿射子集（其坐标记为 $x$、$y$、$z$）中由 $z$-轴给出。（我们使用仿射坐标，是为了帮助看清正在做的事情；尽管如此，这个论证用射影坐标来做会更好。在重读时，你应当把它翻译成一个完全射影的论证。）

\textbf{27.2.C.（练习）}\quad 考虑形如 $(x,y,z)=(a,b,0)+t(a',b',1)$ 的直线（其中 $(a,b,a',b')\in\mathbb{A}^{4}$ 固定，而 $t$ 在 $\mathbb{A}^{1}$ 中变动）。证明：$a$、$b$、$a'$、$b'$ 可以被释读为\trGrassmannian{}的某个标准\trOpenSubset{}上的「通常」坐标（见 \S7.7），并以 $[\ell]$ 为原点。

在\trModuliSpace{}上建立好局部坐标之后，我们就可以进入正题了。设 $f(x,y,z)=0$ 是三次曲面 $X$ 的方程（的仿射版本）。由于 $X$ 包含 $z$-轴 $\ell$，故 $f(x,y,z)\in(x,y)$。更一般地，直线
\[
\tag{27.2.3.1}\label{eq:nearby-line-param}
(x,y,z)=(a,b,0)+t(a',b',1)
\]
含于 $X$，当且仅当 $f(a+ta',\,b+tb',\,t)$ 作为 $t$ 的三次多项式恒为 $0$。这等价于 $a$、$a'$、$b$、$b'$ 的四个方程，分别对应于 $t^{3}$、$t^{2}$、$t$ 与 $1$ 的系数。这比一个单纯的集合论陈述要强：

\textbf{27.2.D.（练习）}\quad 验证：这四个方程是\trSchemeTheoreticFiber（或\trSchemeTheoreticPreimage）$\pi^{-1}([X])$ 的局部方程。

现在我们来到论证的关口，在这里我们要用到 $X$（沿 $\ell$）的正则性。我们有的是一个代数中的精确问题。给定一张由 $f(x,y,z)=0$ 给出的三次曲面 $X$，它含有 $\ell$，而且我们知道 $X$ 是正则的（包括「在无穷远处」，即在 $\mathbb{P}^{3}$ 中）。为了证明 $[\ell]=V(a,a',b,b')$ 是 $\pi$ 的纤维中的一个\trReduced{}孤立点，我们在环 $k[a,a',b,b']/(a,a',b,b')^{2}$ 中工作，也就是说，我们施加方程

\begin{figure}[htbp]
\centering
\begin{tikzpicture}[scale=1.0,line width=0.7pt]
  \draw (0.6,3.4) -- (4.6,3.7) -- (6.2,2.9) -- (2.2,2.6) -- cycle;
  \node[anchor=west] at (5.35,2.72) {$z=1$};
  \draw (0.2,1.0) -- (4.2,1.3) -- (5.8,0.5) -- (1.8,0.2) -- cycle;
  \node[anchor=west] at (4.85,0.32) {$z=0$};
  \fill (2.6,3.06) circle (1.2pt);
  \node[anchor=south east] at (2.62,3.08) {$(0,0,1)$};
  \fill (4.4,3.2) circle (1.2pt);
  \node[anchor=south west] at (4.42,3.22) {$(a+a',\,b+b',\,1)$};
  \fill (2.2,0.86) circle (1.2pt);
  \node[anchor=north east] at (2.22,0.84) {$(0,0,0)$};
  \fill (3.4,0.95) circle (1.2pt);
  \node[anchor=north west] at (3.42,0.93) {$(a,b,0)$};
  \draw (2.6,3.06) -- (2.2,0.86);
  \node[anchor=east] at (2.1,2.0) {$\ell$};
  \draw (4.4,3.2) -- (3.4,0.95);
  \node[anchor=west] at (4.15,2.1) {「附近的直线」};
\end{tikzpicture}
\caption{「$\ell$ 附近的直线」所成空间的参数：以它们与 $z=0$ 平面和 $z=1$ 平面相交的位置表出。（依原书手绘图重绘。）}\label{fig:lines-near-ell}
\end{figure}

\[
\tag{27.2.3.2}\label{eq:nearby-line-square-zero}
a^{2}=aa'=\cdots=(b')^{2}=0,
\]
并试图由此证明 $a=a'=b=b'=0$。（理解「我们为何令 $(a,a',b,b')^{2}=0$」是至关重要的。你也可以从所涉函数的导数这一角度来解读这个论证——归根结底，这可以理解为忘掉高阶信息、只记住相关变量中的一次项；参看\mbox{练习~13.1.I}。关于以导数语言对这一计算的描述，见 [Mu6, \S8D]。）

设 $f(x,y,z)=c_{x^{3}}x^{3}+c_{x^{2}y}x^{2}y+\cdots+c_{1}1=0$，其中 $c_{x^{3}},\,c_{x^{2}y},\,\cdots\in\bar{k}$。由于 $\ell\subset X$，即 $f\in(x,y)$，我们有 $c_{1}=c_{z}=c_{z^{2}}=c_{z^{3}}=0$。现在把 \eqref{eq:nearby-line-param} 代入 $f$，再施加 \eqref{eq:nearby-line-square-zero}。$f$ 的系数中，只有那些恰含一个 $x$ 或一个 $y$ 的单项式的系数存活下来：
\begin{multline*}
c_{x}(a+a't)+c_{xz}(a+a't)t+c_{xz^{2}}(a+a't)t^{2} \\
+c_{y}(b+b't)+c_{yz}(b+b't)t+c_{yz^{2}}(b+b't)t^{2} \\
=\;(a+a't)\bigl(c_{x}+c_{xz}t+c_{xz^{2}}t^{2}\bigr)+(b+b't)\bigl(c_{y}+c_{yz}t+c_{yz^{2}}t^{2}\bigr)
\end{multline*}
被要求作为 $t$ 的多项式恒为 $0$。（回想 $c_{x},\dots,c_{yz^{2}}$ 是 $\bar{k}$ 中固定的元素。）为方便计，令 $C_{x}(t)=c_{x}+c_{xz}t+c_{xz^{2}}t^{2}$ 与 $C_{y}(t)=c_{y}+c_{yz}t+c_{yz^{2}}t^{2}$。

现在，$X$ 在 $(0,0,0)$ 处正则，当且仅当 $c_{x}$ 与 $c_{y}$ 不全为零（因为 $c_{z}=0$）。更一般地，$X$ 在 $(0,0,t_{0})$ 处正则，当且仅当 $c_{x}+c_{xz}t_{0}+c_{xz^{2}}t_{0}^{2}=C_{x}(t_{0})$ 与 $c_{y}+c_{yz}t_{0}+c_{yz^{2}}t_{0}^{2}=C_{y}(t_{0})$ 不全为零。你应当能够很快检验：$X$ 在 $\ell$ 「在无穷远处」的那一点处正则，当且仅当 $c_{xz^{2}}$ 与 $c_{yz^{2}}$ 不全为零。我们把这一点总结如下：$X$ 在 $\ell$ 的每一点处都正则，当且仅当两个二次多项式 $C_{x}(t)$ 与 $C_{y}(t)$ 没有公共根——包括「在无穷远处」。

现在我们就用它，在 $(a+a't)C_{x}(t)+(b+b't)C_{y}(t)\equiv 0$ 的条件下迫使 $a=a'=b=b'=0$。

我们先处理 $C_{x}$ 与 $C_{y}$ 有两个相异根、且两根都有限的特殊情形（即 $c_{xz^{2}}$ 与 $c_{yz^{2}}$ 非零）。设 $t_{0}$ 与 $t_{1}$ 是 $C_{x}(t)$ 的根；把 $t_{0}$ 与 $t_{1}$ 代入 $(a+a't)C_{x}(t)+(b+b't)C_{y}(t)$，便得到 $b+b't_{0}=0$ 与 $b+b't_{1}=0$，从而 $b=b'=0$。同理，$a=a'=0$。

\textbf{27.2.E.（练习）}\quad 处理余下的情形，以完成\mbox{定理~27.2.1} 的证明。（这件事可以做得相当巧妙。例如，你或许能够重新选取坐标，使得 $C_{x}$ 与 $C_{y}$ 的根都是有限的。）\hfill$\square$

\textbf{27.2.4. 直线的构型。}\quad 所谓三次曲面上的「直线构型」，指的是「27 条线中哪些对相交」这份数据。在\trFermat{}三次曲面这一特殊情形（\mbox{练习~27.1.C}），我们不难把它算出来。（利用 $X$ 作为 $\mathbb{P}^{2}$ 的\trBlowup{}的描述，或许更具启发性；见\mbox{练习~27.3.F}。）我们现在来证明：对全体光滑三次曲面，这个构型都是「相同的」（在恰当的解读之下）。

\textbf{27.2.F.（练习）}\quad 构造一个次数为 $27!$ 的有限平展映射 $W\to\mathbb{P}^{19}\setminus\Delta$，它把一张三次曲面连同 27 条相异直线的一个\emph{有序清单}一同参数化。提示：令 $W'$ 为 $Z$ 在 $\mathbb{P}^{19}\setminus\Delta$ 上的第 $27$ 个\textbf{纤维幂}（fibered power），把它释读为对「一张三次曲面连同 27 条直线的有序清单（诸线不必相异）」的参数化。令 $W$ 为诸线相异处对应的子集，证明 $W$ 在 $W'$ 中既开又闭，从而是 $W'$ 的若干连通分量之并。另一种在 $W'$ 中辨认出 $W$ 的办法（这种做法或许能推广到\mbox{练习~27.2.H}）：给「对每一对 $i,j$，『$\ell_i\cdot\ell_j$』是 $W'$ 上的局部常值函数」这一事实赋予意义，于是这些函数对一切 $i\neq j$ 都取正值的位置既开又闭；证明这正是我们所要的 $W$。

现在我们来把下面这个陈述说清楚：\trFermat{}三次曲面（记作 $X_{0}$）上的直线构型与另外某张光滑三次曲面（记作 $X_{1}$）上的构型是「相同的」。把点 $[X_{0}]$ 提升为一点 $w_{0}\in W$。令 $W''$ 为 $W$ 的含 $w_{0}$ 的连通分量。

\textbf{27.2.G.（练习）}\quad 证明：态射 $W''\to\mathbb{P}^{19}\setminus\Delta$ 有限且平展（并且是满的）。

选取映射到 $[X_{1}]$ 的一点 $w_{1}\in W''$。由于 $W$ 参数化的是曲面上 27 条线的「标号」（labeling）即排序，我们由此便选定了一种把 $X_{0}$ 上的诸线与 $X_{1}$ 上的诸线等同起来的方式。设 $X_{0}$ 上的诸线为 $\ell_{1},\dots,\ell_{27}$，而 $X_{1}$ 上相应的诸线为 $m_{1},\dots,m_{27}$。

\textbf{27.2.H.（练习，用到带星的\mbox{练习~24.7.D}）}\quad 证明：对一切 $i$ 与 $j$ 都有 $\ell_{i}\cdot\ell_{j}=m_{i}\cdot m_{j}$。

\textbf{27.2.I.（练习）}\quad 证明：对每张光滑三次曲面 $X\subset\mathbb{P}^{3}_{\bar{k}}$，$X$ 的每条直线都恰好与 $X$ 上另外 $10$ 条直线 $\ell_{1},\ell_{1}',\dots,\ell_{5},\ell_{5}'$ 相交，其中对每个 $i$，$\ell_{i}$ 与 $\ell_{i}'$ 相交，而这些直线之间再无别的相交对。

\textbf{27.2.J.（练习）}\quad 证明：每张光滑三次曲面都含有两条不相交的直线 $\ell$ 与 $\ell'$，使得恰有另外 $5$ 条直线 $\ell_{1},\dots,\ell_{5}$ 与 $\ell$ 和 $\ell'$ 都相交\footnote{原书此处作「precise five other lines」，「precise」当为「precisely」之误，径改。——译者注}。（事实上，对光滑三次曲面中任意两条不相交的直线 $\ell$ 与 $\ell'$，都恰有五条别的直线与二者相交；但这比题目所要求的严格更难。）

\textbf{27.2.5.}\ \emph{\trhRemark（参看\mbox{注记~27.3.7}）：\trWeyl{}群 $W(E_{6})$。}\quad 直线构型的对称群——也就是说，27 条线的置换中保持相交数据者所构成的子群——神奇地恰好是 $E_{6}$ 的\trWeyl{}群，一个阶为 $51840$ 的群。（利用\mbox{练习~27.1.C} 的\trFermat{}三次曲面，以你现有的知识至少足以验证这个群的阶确为 $51840$，不过这要费些功夫。）$W''$ 在 $\mathbb{P}^{19}\setminus\Delta$ 上的次数是 $51840$，而 $K(Z)/K(\mathbb{P}^{19}\setminus\Delta)$ 的\trGalois{}闭包的\trGalois{}群同构于 $W(E_{6})$（见 [H1, III.3]）——这绝非巧合。

\textbf{27.2.K.}\ \extrahard\ \textbf{（练习）}\quad 在任意特征下证明\mbox{定理~27.2.1}。先把\mbox{练习~27.2.A} 在 $\mathbb{Z}$ 上的正确陈述想明白并加以证明。然后遵循本节所给的论证，在必要时作出改动。

\textbf{27.2.6.}\ \hard\ \textbf{直线的 Fano 簇与\trHilbertScheme。}\quad 在\mbox{练习~27.2.A} 与\mbox{练习~27.2.K} 中，你已经把含于某张 $X$ 中的直线的\trModuliSpace{}作为一个概形构造了出来。\footnote{原书此处作「contained in a X」，当为「contained in an X」之误，径改。——译者注}你的论证可以推广到任何 $X\subset\mathbb{P}^{N}$。这一构造称为 $X$ 的\textbf{直线的 Fano 簇}（Fano variety of lines）（请勿与 \S21.5.5 中的 Fano 簇（Fano variety）概念相混淆），它是\trHilbertScheme{}的一个例子。
